

\begin{derivation}[从\eqnrefs{eq:plemelj-split-r13}到\eqnrefs{eq:sokhotski-plemelj-dp68}]

正文\eqnrefs{eq:plemelj-split-r13} 已把复分母分成实部和虚部两个普通函数。$\epsilon\to0^+$ 时，实部按正文\eqnrefs{eq:principal-value-definition-dp68} 趋于 Cauchy 主值，虚部则形成单位面积、宽度趋零的 Lorentz 峰；分别取这两个分布极限即可得到正文\eqnrefs{eq:sokhotski-plemelj-dp68}。

先令
\begin{equation*}
 \delta_\epsilon(x)
 =\frac1\pi\frac{\epsilon}{x^2+\epsilon^2}.
\end{equation*}
变量代换 $x=\epsilon u$ 给
\begin{equation*}
 \int\dd x\,\delta_\epsilon(x)=1,
\end{equation*}
且对任意光滑测试函数 $f$，
\begin{equation*}
 \int\dd x\,\delta_\epsilon(x)f(x)
 =\frac1\pi\int\frac{\dd u}{1+u^2}f(\epsilon u)
 \longrightarrow f(0).
\end{equation*}
所以 $\delta_\epsilon\to\delta$。另一方面，把
$x/(x^2+\epsilon^2)$ 的正负半轴配对后，极限正是正文\eqnrefs{eq:principal-value-definition-dp68}。代回正文\eqnrefs{eq:plemelj-split-r13} 即得正文\eqnrefs{eq:sokhotski-plemelj-dp68}。
\end{derivation}

\begin{derivation}[从\eqnrefs{eq:jordan-exponential-r13}到\eqnrefs{eq:contour-orientation-residue-dp68}]

正文\eqnrefs{eq:jordan-exponential-r13} 已告诉我们指数在复平面的哪一侧衰减。还需要再检查大半圆积分确实消失以及闭合曲线的取向：$\tau>0$ 时下半圆顺时针，$\tau<0$ 时上半圆逆时针。两点合起来便得到正文\eqnrefs{eq:contour-orientation-residue-dp68} 的正负号。

设大 $|z|$ 时 $|F(z)|\le C/|z|^{1+\delta}$，$\delta>0$。对 $\tau>0$ 取下半圆 $z=Re^{-\ii\theta}$，则
\begin{equation*}
 |\ee^{-\ii z\tau/\hbar}|=\ee^{-R\tau\sin\theta/\hbar}\le1,
\end{equation*}
从而
\begin{equation*}
 \left|\int_{C_R^-}\dd z\,F(z)\ee^{-\ii z\tau/\hbar}\right|
 \le\frac{\pi C}{R^\delta}\to0.
\end{equation*}
下半圆从 $+R$ 绕回 $-R$，闭合曲线顺时针，因此留数定理给实轴积分
$-2\pi\ii$ 乘下半平面留数；除以正文 Fourier convention 中的 $2\pi$ 即得\eqnrefs{eq:contour-orientation-residue-dp68} 第一行。$\tau<0$ 时同理向上闭合，曲线逆时针，得到第二行。
\end{derivation}

